\section{Student inference: complete proofs}
\label{sec:annex-p15-student-inference}

\begin{annexitem}{annex:p15.student.orthogonal}{thm:p15-student-orthogonal}{The orthogonal construction and chi-square premise}
\begin{proof}
For \(n\geq2\), the vector \(v=b-d\) is nonzero and
\(v^{\mathsf T}v=2(1-n^{-1/2})\).
The matrix \(P=vv^{\mathsf T}/(v^{\mathsf T}v)\) satisfies
\(P^{\mathsf T}=P=P^2\). Therefore \(O=I-2P\) is symmetric
and \(O^{\mathsf T}O=I\). Since
\(v^{\mathsf T}b=1-n^{-1/2}\), we have
\(Ob=b-v=d\); symmetry gives \(b^{\mathsf T}O=d^{\mathsf T}\).
This proves the claimed first row.

The joint density of \((X-\mu\mathbf1)/\sigma\) is
\((2\pi)^{-n/2}\exp(-\|y\|^2/2)\), by independence and
one-dimensional translation and scaling.
\Cref{thm:p15-wce-spherical-area} proves that an orthogonal
map preserves Lebesgue measure. Its use here requires only
the already verified orthogonality of \(O\).
Since \(O\) also preserves the squared norm, the transformed
vector has the same density. This density is the product of
the \(n\) normalized standard normal densities
of \cref{thm:p15-normal-normalization}; integrating over
rectangles proves independence of its coordinates.
The first row gives
\(Y_1=\sqrt n(\bar X-\mu)/\sigma\). The identity
\[
 \sum_i(X_i-\mu)^2
       =\sum_i(X_i-\bar X)^2+n(\bar X-\mu)^2
\]
follows by expansion, since \(\sum_i(X_i-\bar X)=0\).
Subtracting \(Y_1^2\) from the squared norm of \(Y\)
therefore gives \(W=V/\sigma^2\).

If \(G\) is standard normal, differentiation of
\(P(G^2\leq q)=\int_{-\sqrt q}^{\sqrt q}
 (2\pi)^{-1/2}e^{-y^2/2}\,dy\), \(q>0\), gives density
\[
                       \frac{q^{-1/2}e^{-q/2}}{\sqrt{2\pi}}.
\]
The gamma integral, after \(q=u^2\), gives
\(\Gamma(1/2)=2\int_0^\infty e^{-u^2}\,du=\sqrt\pi\)
by the proved Gaussian integral.
Thus \(G^2\) is \(\mathrm{Gamma}(1/2,1/2)\).
The \(\nu\) squares in \(W\) are independent, so repeated
application of \cref{thm:p15-beta-gamma-convolution}
gives \(\mathrm{Gamma}(\nu/2,1/2)\); for \(\nu=1\)
no convolution is needed.
This gamma density has total mass one on \(q>0\)
by \cref{thm:p15-gamma-normalization}. Hence \(W>0\)
almost surely. Finally, \(Y_1\) is independent of the
vector \((Y_2,\ldots,Y_n)\), and thus of its measurable
function \(W\).
\end{proof}
\end{annexitem}

\begin{annexitem}{annex:p15.student.beta.direction}{thm:p15-student-beta-direction}{The ratio calculation and the sphere coordinate}
\begin{proof}
Let \(A,B\) have the densities in the statement.
For any nonnegative Borel function \(H(q,z)\), their joint
integral is
\[
 \frac{\lambda^{a+b}}{\Gamma(a)\Gamma(b)}
 \int_0^\infty\int_0^\infty
 H\!\left(u+v,\frac{u}{u+v}\right)
 u^{a-1}v^{b-1}e^{-\lambda(u+v)}\,dv\,du .
\]
For fixed \(u\), substitute \(q=u+v\). Tonelli then
exchanges the region \(0<u<q<\infty\), and for each fixed
\(q\) the one-dimensional substitution \(u=qz\) gives
\[
 \int_0^\infty\int_0^1 H(q,z)
 \frac{\lambda^{a+b}q^{a+b-1}e^{-\lambda q}}
      {\Gamma(a+b)}
 \frac{z^{a-1}(1-z)^{b-1}}{B(a,b)}\,dz\,dq .
\]
The beta-gamma identity
\(B(a,b)=\Gamma(a)\Gamma(b)/\Gamma(a+b)\)
is exactly \cref{thm:p15-beta-gamma-convolution}.
Both displayed factors are normalized densities, proving
the two laws and independence by taking rectangle indicators.
This derivation supplies the ratio law rather than assuming
it from the gamma sum law.

For a standard normal \(Y_1\), symmetry of its density
shows that its sign is fair and independent of \(Y_1^2\):
the probabilities of \(Y_1>0,Y_1^2\in D\) and
\(Y_1<0,Y_1^2\in D\) are both half
\(P(Y_1^2\in D)\). A zero value has probability zero.
The sign is also independent of the remaining normal
coordinates, so apply the ratio result to
\[
       A=Y_1^2,\qquad B=\sum_{j=2}^nY_j^2
\]
with shapes \(1/2,\nu/2\) and rate \(1/2\).
It follows that \(Y_1^2/\|Y\|^2\) is
\(\mathrm{Beta}(1/2,\nu/2)\), independently of
\(\|Y\|^2\), and its signed square root has density
\[
 \frac{\Gamma(n/2)}{\Gamma(1/2)\Gamma(\nu/2)}
 (1-z^2)^{\nu/2-1},\qquad -1<z<1.
\]
Indeed on either half-axis the change from \(r=z^2\)
contributes \(2|z|\), while the fair sign contributes
\(1/2\); these cancel the \(r^{-1/2}\) factor.
This proves the displayed \(h_n\), including its
normalization, symmetry and absence of atoms.

The remaining vector assertion is the radius/direction
conclusion of \cref{thm:p15-wce-sphere-coordinate}.
Its hypotheses are exactly independent standard normal
coordinates, as just used; its direction measure is the
same normalized cone measure \(\omega_n=\sigma_n\).
That result follows from the polar formula in
\cref{thm:p15-wce-spherical-area}, so no conditional
distribution on a set of positive radial thickness is
being confused with surface measure.
Orthogonal invariance of \(\omega_n\) and the matrix
constructed in \cref{thm:p15-student-orthogonal}
show that \(d^{\mathsf T}U\) has the same law as
the first coordinate \(Y_1/\|Y\|\). The law of
\(\|Y\|^2\) is the sum of \(n\) independent squared
normals, hence \(\chi_n^2\), as above.
\end{proof}
\end{annexitem}

\begin{annexitem}{annex:p15.student.law}{thm:p15-student-law}{The full normalized density}
\begin{proof}
Condition on \(W=w>0\). The scaled normal
\(Z_0/\sqrt{w/\nu}\) has density
\(\sqrt{w/\nu}(2\pi)^{-1/2}
 \exp(-t^2w/(2\nu))\).
Multiplying by the gamma density of \(W\) and integrating
over \(w>0\) gives
\[
\begin{split}
 f_\nu(t)
 &=\frac{1}{\sqrt{2\pi\nu}\,2^{\nu/2}\Gamma(\nu/2)}
   \int_0^\infty
       w^{(\nu+1)/2-1}
       e^{-(1+t^2/\nu)w/2}\,dw\\
 &=\frac{\Gamma((\nu+1)/2)}
          {\sqrt{\nu\pi}\,\Gamma(\nu/2)}
       (1+t^2/\nu)^{-(\nu+1)/2}.
\end{split}
\]
The second equality is the gamma rate substitution
already proved in \cref{thm:p15-gamma-normalization}.
Tonelli and integration of the conditional normal
density show \(\int_{\mathbb R}f_\nu(t)\,dt=1\);
thus the calculation establishes normalization as well
as the functional form.
The density is positive, continuous and even. Its
integrals therefore give a continuous strictly
increasing CDF from zero to one, with the stated
reflection identity and unique finite quantiles.

The first theorem gives
\[
 T_\mu=\frac{\sqrt n(\bar X-\mu)/\sigma}
                 {\sqrt{(V/\sigma^2)/\nu}},
\]
with independent standard normal numerator and
\(\chi_\nu^2\) denominator square. This proves the
nuisance-free pivot law.
For \(Z=Y_1/\sqrt{Y_1^2+W}\), direct algebra gives
\(\sqrt\nu Z/\sqrt{1-Z^2}=Y_1/\sqrt{W/\nu}\).
The denominator is positive almost surely.
By the preceding spherical-coordinate proof, this
is the law asserted for the abstract coordinate \(Z\).
\end{proof}
\end{annexitem}

\begin{annexitem}{annex:p15.student.coverage}{thm:p15-student-coverage}{Inversion, both tails, and zero residuals}
\begin{proof}
On \(V>0\),
\[
\begin{aligned}
 \mu\in C_-(X)&\iff T_\mu\leq t_{\nu,\alpha},\\
 \mu\in C_+(X)&\iff T_\mu\geq-t_{\nu,\alpha},\\
 \mu\in C_0(X)&\iff
            |T_\mu|\leq t_{\nu,(1+\alpha)/2}.
\end{aligned}
\]
The first two probabilities are \(\alpha\), since
the CDF is continuous and symmetric. In the third,
the critical value is positive and the probability is
\[
 F_\nu(t_{\nu,(1+\alpha)/2})
       -F_\nu(-t_{\nu,(1+\alpha)/2})
                       =\alpha.
\]
These computations hold for every true \(\mu\) and
every positive \(\sigma\). The common conclusion
\(P_{\mu,\sigma}(V=0)=0\) shows that the convention
there cannot change coverage.

For every sample vector, including \(V=0\), solving
the complement of each displayed strict rejection
inequality for the candidate \(m\) gives precisely
the three closed intervals in the statement.
At \(e=0\) their endpoints reduce to \(\bar x\);
no comparison involving the undefined ratio \(T_m\)
is used. The Borel character of all endpoints and
inclusion relations follows from sums, square roots
of nonnegative continuous functions, and inequalities.
Finally the one-sided table tail is
\(P(|T_\mu|>t_{\nu,\alpha})=2(1-\alpha)\) when
\(\alpha>1/2\); the central table tail is
\(1-\alpha\). These are the two distinct instructions
in the source. Signed quantiles give the stated
one-sided formulas throughout \(0<\alpha<1\).
\end{proof}
\end{annexitem}

\begin{annexitem}{annex:p15.student.laplace.completeness}{thm:p15-student-laplace-completeness}{A signed Laplace uniqueness proof}
\begin{proof}
Fix \(\lambda_0>0\). The hypothesis makes
\[
        \eta(dq)=g(q)q^{a-1}e^{-\lambda_0q}\,dq
\]
a finite signed measure on \((0,\infty)\).
Push it forward by the Borel bijection
\(q\mapsto z=e^{-q}\) onto \((0,1)\), and extend
the resulting signed measure \(\rho\) to \([0,1]\)
by assigning zero mass to the endpoints.
For every integer \(k\geq0\),
\[
       \int_{[0,1]} z^k\,d\rho(z)
       =\int_0^\infty g(q)q^{a-1}
                    e^{-(\lambda_0+k)q}\,dq=0.
\]
Thus all polynomial integrals vanish.
\Cref{thm:p15-feller-bernstein} supplies, for every
continuous \(f\) on \([0,1]\), polynomials converging
uniformly to \(f\). Since
\[
 \left|\int(f-p)\,d\rho\right|
      \leq\|f-p\|_\infty\,|\rho|([0,1]),
\]
all continuous integrals vanish as well.

For \(0\leq t<1\), choose continuous functions equal
to one on \([0,t]\), decreasing linearly to zero on
\([t,\min(1,t+1/j)]\), and zero farther to the right,
using \(j\) large enough that \(t+1/j\leq1\).
They are bounded by one and converge pointwise to
\(\mathbf1_{[0,t]}\). Dominated convergence with
respect to the finite total variation \(|\rho|\)
gives \(\rho([0,t])=0\). The case \(t=1\)
is the zero integral of the constant polynomial.
These initial intervals form a generating
intersection-closed family \(\mathcal P\) of Borel sets.
Because \(\rho([0,1])=0\), its zero-measure sets form
a class \(\mathcal D\) containing the whole space and
closed under complements and countable disjoint unions.
Here is the required intersection-class argument explicitly.
Let \(\mathcal L\) be the smallest class with these
three closure properties that contains \(\mathcal P\);
then \(\mathcal L\subseteq\mathcal D\).
Such a class is closed under \(B\setminus A\)
whenever \(A\subseteq B\) and both belong to the class:
take the complement of the disjoint union \(A\cup B^c\).
For fixed \(A\in\mathcal P\), the sets \(B\) such that
\(A\cap B\in\mathcal L\) have the three closure properties,
by this relative-complement observation, and contain
\(\mathcal P\). They therefore contain \(\mathcal L\).
Now fix \(B\in\mathcal L\) and repeat the same argument
with the roles interchanged. This proves that
\(\mathcal L\) is closed under intersections.
It is consequently closed under finite unions and,
by disjointizing any sequence using finite unions
and complements, under arbitrary countable unions.
Thus \(\mathcal L\) is a sigma algebra containing
\(\mathcal P\), so contains all Borel sets.
All those sets lie in \(\mathcal D\), proving \(\rho=0\).
The Borel bijection with Borel inverse then gives
\(\eta=0\).

The density of \(\eta\) consequently vanishes almost
everywhere: otherwise a set where it is at least
\(1/j\), or at most \(-1/j\), intersected with a
bounded interval, would have a nonzero integral.
Since \(q^{a-1}e^{-\lambda_0q}>0\) on \((0,\infty)\),
this proves \(g=0\) almost everywhere.
Multiplication by the positive gamma normalizing
constant turns the final assertion exactly into
the displayed Laplace premise.
\end{proof}
\end{annexitem}

\begin{annexitem}{annex:p15.student.null.fibers}{thm:p15-student-null-fibers}{Every radius law and every alternative direction}
\begin{proof}
Under \(P_{m,\sigma}\),
\((X-m\mathbf1)/\sigma\) has independent standard
normal coordinates. Its normalized direction is
\(\omega_n\), independently of its squared radius,
by \cref{thm:p15-student-beta-direction}.
Scaling \(\chi_n^2\) by \(\sigma^2\) gives density
\[
 g_\sigma(q)=
 \frac{q^{n/2-1}e^{-q/(2\sigma^2)}}
           {(2\sigma^2)^{n/2}\Gamma(n/2)},\qquad q>0.
\]
For any nonnegative Borel \(h\), the resulting
integral identity is
\[
 E_{m,\sigma}h(X)
   =\int_0^\infty g_\sigma(q)
       \int h(m\mathbf1+\sqrt q\,u)\,d\omega_n(u)\,dq.
\]
This explicitly supplies the conditional kernel,
including its measurability, without a choice of
versions depending on \(\sigma\).

At the same positive variance, expansion of the
normal joint densities gives
\[
 \frac{p_{m+\delta,\sigma}(x)}{p_{m,\sigma}(x)}
       =\exp\!\left(
          \frac{\delta\sum_i(x_i-m)}{\sigma^2}
                       -\frac{n\delta^2}{2\sigma^2}\right).
\]
At \(x=m\mathbf1+\sqrt q\,u\), its first exponent
is \(a(q)d^{\mathsf T}u\).
Insert this ratio in the preceding integral
identity. Integrating over \(u\) gives
\(g_\sigma(q)e^{-n\delta^2/(2\sigma^2)}M(a(q))\);
normalizing the inner integral gives
\(e^{a(q)d^{\mathsf T}u}/M(a(q))\).
For finite \(a\), \(e^{-|a|}\leq M(a)\leq e^{|a|}\);
therefore the kernels are finite and positive.
Their total integral is one because they reproduce
the normalized alternative normal density.
This proves every asserted law at both positive
and negative \(\delta\), and also at \(\delta=0\).

Finally
\[
 Q_m=V+n(\bar X-m)^2,\qquad
 Z_m^2=\frac{n(\bar X-m)^2}{Q_m}
\]
by the same centered-sum identity used earlier.
On \(V>0\), these equations give \(|Z_m|<1\)
and the displayed signed relation to \(T_m\).
The cases \(Q_m=0\) and \(V=0\) are null under
all the normal laws; the definitions on them are
already specified and no density there is needed.
\end{proof}
\end{annexitem}

\begin{annexitem}{annex:p15.student.score.fibers}{thm:p15-student-score-fibers}{Differentiation and the complete nuisance constraints}
\begin{proof}
Fix \(\sigma>0\) and write \(y=x-m\mathbf1\).
For \(|\delta|\leq D\),
\[
  \|y-\delta\mathbf1\|^2
             \geq \tfrac12\|y\|^2-nD^2.
\]
The first two mean derivatives of the normal
density are that density multiplied respectively
by \(\sum_i(x_i-\mu)/\sigma^2\) and by
\((\sum_i(x_i-\mu))^2/\sigma^4-n/\sigma^2\).
On this neighborhood they are bounded in absolute
value by a constant times
\((1+\|y\|^2)e^{-\|y\|^2/(4\sigma^2)}\).
This function is integrable by the normal
normalization and moment formulas; \(|\phi|\leq1\)
does not change the envelope.
Differentiation under the integral and dominated
convergence give continuous first and second
derivatives and the two stated score identities.
The same argument gives continuity of power itself.

The angular integrands defining \(A_\phi,B_\phi\)
are Borel, so their integrals are Borel in \(q\).
They obey \(0\leq A_\phi\leq1\), \(|B_\phi|\leq1\).
The null-fiber formula gives
\[
        E_{m,\sigma}\phi
                         =\int_0^\infty A_\phi(q)g_\sigma(q)\,dq.
\]
If this equals \(\varepsilon\) for every
\(\sigma>0\), the rates \(1/(2\sigma^2)\) range
over all of \((0,\infty)\).
\Cref{thm:p15-student-laplace-completeness},
with \(a=n/2\) and bounded function
\(A_\phi-\varepsilon\), gives
\(A_\phi=\varepsilon\) Lebesgue almost everywhere.
The converse follows immediately by integration
against \(g_\sigma\).

On a sphere \(H_m=\sqrt{nq}\,d^{\mathsf T}u\);
the first-derivative formula is consequently
\[
       \partial_\mu\beta_\phi(m,\sigma)
        =\frac{\sqrt n}{\sigma^2}
              \int_0^\infty \sqrt q\,B_\phi(q)g_\sigma(q)\,dq.
\]
The function \(\sqrt q\,B_\phi(q)\) is integrable
against every gamma density here, since it is
bounded by \(\sqrt q\) and gamma has that moment
by its defining integral.
Vanishing of all these derivatives and the
same completeness theorem yield \(B_\phi=0\)
almost everywhere on \(q>0\).
The converse again follows by integration.
Taking the union of these two Lebesgue null sets
supplies one null set sufficient for all values
of the nuisance parameter. Every alternative
radius density in the preceding theorem is
absolutely continuous with respect to Lebesgue
measure, so it also assigns this set zero mass.

If \(\phi\) is globally unbiased of level
\(\varepsilon\), then at each fixed \(\sigma\)
continuity in \(\mu\) and the lower bound
\(\beta_\phi(\mu,\sigma)\geq\varepsilon\) for
all \(\mu\ne m\) imply
\(\beta_\phi(m,\sigma)\geq\varepsilon\).
The level condition gives the reverse inequality.
Thus equality holds, and this value is a global
minimum of the differentiable power as a
function of \(\mu\). Its first derivative is
zero. Both conclusions hold for every
\(\sigma>0\), which is stationary similarity.
To see why the angular conclusion is only almost everywhere,
start with the similar upper-cap test and set it equal to
zero on the single sphere \(Q_m=q_0\), where \(q_0>0\).
This Borel sphere has probability zero under every normal
law in the model, so the altered test is still similar.
Its angular mass at \(q_0\) is zero, whereas
\(\varepsilon>0\). Thus a conclusion for literally every
radius would be false for the stated class of Borel tests.
\end{proof}
\end{annexitem}

\begin{annexitem}{annex:p15.student.one.sided.optimal}{thm:p15-student-one-sided-optimal}{The two orientations of the conditional comparison}
\begin{proof}
The strictly increasing map
\(z\mapsto\sqrt\nu z/\sqrt{1-z^2}\) sends
\((-1,1)\) onto \(\mathbb R\), with inverse
\(t\mapsto t/\sqrt{\nu+t^2}\).
The pivot and spherical laws show
\(\omega_n(Z>c_+)=1-\alpha=\varepsilon\).
Thus the upper test is, up to its null
zero-residual set, the same upper cap on
every null sphere and is similar.

Let \(\phi\) be any similar test.
For almost every \(q>0\), put
\(g(u)=\phi(m\mathbf1+\sqrt q\,u)\).
Then \(\int g\,d\omega_n=\varepsilon\).
At an alternative \(\mu=m+\delta>m\),
\(a(q)>0\). The pointwise inequality
\[
  \bigl(\mathbf1_{\{Z>c_+\}}-g(u)\bigr)
                  \bigl(e^{a(q)Z}-e^{a(q)c_+}\bigr)\geq0
\]
holds: both factors are nonnegative above the
cut and nonpositive below it. Integrate and
use the equal masses to cancel the constant
\(e^{a(q)c_+}\). This proves domination of
the numerator of the conditional alternative
power. Division by the positive \(M(a(q))\)
and integration against the alternative radius
density prove unconditional domination at
every \(\delta>0,\sigma>0\).

For the lower test, the null cap \(Z<-c_+\)
also has mass \(\varepsilon\).
If \(\delta<0\), \(e^{a(q)Z}\) is decreasing.
The same product inequality, now subtracting
\(e^{-a(q)c_+}\) and using the lower-cap indicator,
has both factors nonnegative below \(-c_+\)
and nonpositive above it. It yields the full
reflected result. This proof permits arbitrary
measurable randomized competitors, arbitrary
positive nuisance variance, and positive,
zero or negative \(c_+\). Atomlessness makes
the cap boundaries irrelevant to the
probability calculation; the unscaled tests
retain the exact endpoint conventions.

For a confidence procedure with exact coverage,
fix the candidate \(m\) and let its rejection
function be one minus its inclusion function.
At the true value \(\mu=m\), this is a similar
test of size \(\varepsilon\) for every \(\sigma\).
If the actual mean is \(\mu>m\), upper-test
domination is precisely minimal inclusion
probability of this false \(m\) by \(C_-\).
If \(\mu<m\), the lower-test result gives
the analogous conclusion for \(C_+\).
This argument specifies the side of the
false parameter and the comparison class
required in the original one-sided assertion.
\end{proof}
\end{annexitem}

\begin{annexitem}{annex:p15.student.exponential.chord}{thm:p15-student-exponential-chord}{The affine chord, strict unbiasedness, and curvature}
\begin{proof}
For \(a=0\) the first inequality is equality
by the prescribed masses. Let \(a\ne0\).
The affine function through the two points
\((-c,e^{-ac})\) and \((c,e^{ac})\) is
\[
        \ell_a(z)=\cosh(ac)+\frac{\sinh(ac)}c\,z.
\]
The function \(z\mapsto e^{az}\) is strictly
convex, so \(e^{az}<\ell_a(z)\) on \((-c,c)\).
For \(z>c\), the secant slope from \(c\) to
\(z\) strictly exceeds the secant slope from
\(-c\) to \(c\), giving
\(e^{az}>\ell_a(z)\); the analogous ordered
secant slopes give this same strict inequality
for \(z<-c\). These statements follow directly
by integrating the strictly increasing
derivative \(a e^{az}\), and hold also when
\(a<0\). At the two cut points equality holds.
Therefore
\[
                   (r-g)(e^{aZ}-\ell_a(Z))\geq0
\]
pointwise. Symmetry gives
\(\int rZ\,d\omega_n=0\), while \(\int r=\varepsilon\).
Both the constant and the linear terms vanish
upon integrating \((r-g)\ell_a(Z)\).
This proves \(\int r e^{aZ}\geq\int g e^{aZ}\)
with both constraints accounted for.

For strict unbiasedness, \(r\) is even in
\(Z\), as is the base density. Thus
\[
 \int r e^{aZ}\,d\omega_n=\int r\cosh(aZ)\,d\omega_n,
 \qquad
 \int e^{aZ}\,d\omega_n=\int\cosh(aZ)\,d\omega_n.
\]
For \(a\ne0\), \(\cosh(az)\) is strictly
increasing with \(|z|\).
Consequently
\[
       (r-\varepsilon)
              \bigl(\cosh(aZ)-\cosh(ac)\bigr)>0
\]
except at the cut points: outside, both factors
are positive; inside, both are negative.
The inside and outside sets have positive
probability because \(0<c<1\) and \(h_n\)
is positive on \((-1,1)\).
Integration, using \(\int(r-\varepsilon)=0\),
therefore gives
\(\int r\cosh(aZ)>\varepsilon\int\cosh(aZ)\).
This proves the strict normalized inequality,
independently of the local curvature calculation.

Finally, for any \(0\leq g\leq1\) of the same
mass, the pointwise product
\[
                         (r-g)(Z^2-c^2)\geq0
\]
has the required sign both inside and outside.
After integration the constant \(c^2\) cancels.
No first-moment constraint is needed for
this last quadratic comparison.
\end{proof}
\end{annexitem}

\begin{annexitem}{annex:p15.student.two.sided.optimal}{thm:p15-student-two-sided-optimal}{From conditional inequalities to global and local optima}
\begin{proof}
The spherical/Student relation makes the
central test equal to
\(r(U_m)=\mathbf1_{\{|Z_m|>c\}}\) whenever
\(V>0\), with \(c\) as displayed.
The central quantile gives
\(\int r\,d\omega_n=\varepsilon\), and
symmetry gives \(\int Zr\,d\omega_n=0\).
Thus \cref{thm:p15-student-score-fibers}
proves stationary similarity.

Fix any stationary similar competitor \(\phi\).
For Lebesgue almost every \(q>0\), its angular
function \(g_q(u)=\phi(m\mathbf1+\sqrt q\,u)\)
has mass \(\varepsilon\) and first moment zero.
Apply \cref{thm:p15-student-exponential-chord}
with \(a=a(q)\). For every alternative
\(\delta\ne0\) and every \(q>0\), this is a
nonzero real number. Division by \(M(a(q))\)
gives central-test domination of conditional
alternative power at almost every radius.
The common exceptional set is null for the
alternative radius density, and integration
gives unconditional domination.
The strict inequality against the constant
\(\varepsilon\) holds at every \(q>0\).
Its integral against the probability density
of the alternative radius is strictly positive,
so central-test power is greater than
\(\varepsilon\) at every alternative.
This proves global unbiasedness.
Every globally unbiased level test is
stationary similar by the score-fiber theorem,
so the domination proves UMPU.

For the local criterion, let \(\phi\) be
any similar competitor; its first derivative
need not vanish for this comparison.
The quadratic part of the conditional lemma
gives, for almost every \(q\),
\[
 \int Z^2\{r(u)-g_q(u)\}\,d\omega_n(u)\geq0.
\]
The second-derivative identity and
\(H_m^2=nqZ_m^2\) yield
\[
\begin{split}
 &\partial_\mu^2\beta_{\phi_m^0}(m,\sigma)
        -\partial_\mu^2\beta_\phi(m,\sigma)\\
 &\quad=\frac n{\sigma^4}
      \int_0^\infty q g_\sigma(q)
          \int Z^2\{r-g_q\}\,d\omega_n\,dq
       \ \geq\ 0 .
\end{split}
\]
Here the term \(-n/\sigma^2\) cancels because
both tests have the same null size.
The integrals are finite since \(Z^2\leq1\)
and the gamma radius has a finite first moment.
The central test itself has zero first
derivative, hence lies in the original
stationary class and maximizes its second
derivative for every nuisance value.
Changing from rejection probability \(\beta\)
to acceptance probability \(1-\beta\)
reverses the second-derivative inequality
and leaves the zero first derivative unchanged.
This is exactly the local criterion.
\end{proof}
\end{annexitem}

\begin{annexitem}{annex:p15.student.unbiased.confidence}{thm:p15-student-unbiased-confidence}{The confidence comparison and its quantifiers}
\begin{proof}
Fix \(m\) and define the bounded Borel test
\(\phi_m(x)=1-a(x,m)\). The coverage condition
at true mean \(m\) gives
\(E_{m,\sigma}\phi_m\leq1-\alpha=\varepsilon\).
For every \(\mu\ne m\), the false-coverage
condition gives \(E_{\mu,\sigma}\phi_m\geq\varepsilon\).
Thus it is a globally unbiased level test.
By \cref{thm:p15-student-score-fibers} its
null size is exactly \(\varepsilon\), for
every \(\sigma>0\). This proves exact true
coverage of \(a\) for every \(m,\sigma\).

UMPU domination gives
\(E_{\mu,\sigma}\phi_m^0\geq E_{\mu,\sigma}\phi_m\).
Taking complements is the claimed
false-inclusion inequality. The Student test's
own strict unbiasedness and similarity prove
the confidence interval's unbiasedness and
coverage. The additional local statement is
the acceptance form of the second-derivative
comparison just proved, applied separately
to each fixed candidate \(m\).

All derivatives in this reasoning vary the
true sampling mean while holding the
candidate \(m\), the procedure, and the
positive nuisance variance fixed.
No derivative of a data-dependent endpoint
or of the candidate-indexed family is required.
Joint measurability makes the confidence
interpretation well defined; the testing
argument uses only its Borel section at
each \(m\). If randomized sets are desired,
an independent uniform variable \(R\) realizes
the inclusion function by
\(m\in C(X,R)\) exactly when \(R\leq a(X,m)\);
the same inequalities follow by integrating
over \(R\).
Finally all probability measures here are
the indexed sampling laws \(P_{\mu,\sigma}\).
Neither the statement nor the proof has
introduced a probability measure on the
parameter space. The asserted conclusions
are consequently the stated sampling and
false-coverage comparisons.
\end{proof}
\end{annexitem}

\begin{annexitem}{annex:p15.student.source.score.bridge}{thm:p15-student-source-score-bridge}{Every normal premise of the source construction}
\begin{proof}
For \(v>0\), the log joint density is
\[
   -\frac n2\log(2\pi v)
                      -\frac1{2v}\sum_i(x_i-\mu)^2.
\]
Differentiating in \(\mu,v\), and then
putting \(\mu=m\), gives
\(\varphi_1=H_m/v\) and
\(\varphi_2=-n/(2v)+Q_m/(2v^2)\).
The second derivatives are
\[
 \varphi_{11}=-n/v,\qquad
 \varphi_{12}=-H_m/v^2=-\varphi_1/v,
 \qquad
 \varphi_{22}=n/(2v^2)-Q_m/v^3
                =-n/(2v^2)-2\varphi_2/v.
\]
These are the affine relations required in
the source equations~(15)--(17).
The density is positive everywhere, and
all the derivative integrals used above
are justified by the score-fiber proof.
More generally, on a compact neighborhood
with \(v\) bounded away from zero and
\(\mu\) bounded, each finite-order \(v\)
derivative and each mixed derivative with
at most two \(\mu\) derivatives is a
polynomial in the sample coordinates
times the normal density, with bounded
coefficients. Completing a square gives
an integrable polynomial times a fixed
Gaussian envelope. This proves the
source's additional differentiation
interchanges whenever they are invoked.

The map from \((H_m,Q_m)\) to
\((\varphi_1,\varphi_2)\) is an invertible
affine map at each fixed \(v>0\).
For sample coordinates \(x_1,x_2\), the
Jacobian of \((H_m,Q_m)\) is
\[
  \det\begin{pmatrix}
    1&1\\2(x_1-m)&2(x_2-m)
  \end{pmatrix}=2(x_2-x_1),
\]
which vanishes only on a Lebesgue null
hyperplane. This verifies the relevant
nondegeneracy even when \(n=2\).
The polar and orthogonal formulas already
proved supply the actual integral
transformation globally; no global
one-to-one property of this local
Jacobian map is assumed.

Fixing \(\varphi_2\) fixes \(Q_m=q\).
The source's conditional preservation
of null mass is \(A_\phi(q)=\varepsilon\);
its conditional first-score equation is
\[
       \int \varphi_1\phi\,d\omega_n
              =\frac{\sqrt{nq}}v B_\phi(q)=0.
\]
These assertions have been proved for
almost every \(q\) through gamma-family
completeness. This supplies precisely
the conditional consequence needed
from the source's more general
orthogonal-polynomial condition.
There is no appeal to unproved
completeness of an arbitrary polynomial
system.

The second power derivative integrates
\(\varphi_{11}+\varphi_1^2\).
Its constant \(\varphi_{11}=-n/v\)
term is fixed by conditional mass,
and its remaining term is
\((nq/v^2)\int Z^2\phi\,d\omega_n\).
The conditional quadratic inequality
selects the two tails \(|Z|>c\);
symmetry makes their first-score
integral zero. This is the actual
normal instance of the source's
outside-two-roots construction and
proves all the premises used for
its local optimum.

For the 1933 parameterization by
\(\sigma\), the nuisance score is
\(-n/\sigma+Q_m/\sigma^3\) and its
derivative is
\(-2n/\sigma^2-3\varphi/\sigma\).
Its level surfaces are the very same
spheres. With the source's
\(s_{33}^2=V/n\), its statistic
\(z_{33}=(\bar X-m)/s_{33}\) equals
\(T_m/\sqrt\nu\). Substitution
\(t=\sqrt\nu z_{33}\) in the normalized
Student density gives density
\[
   \frac{\Gamma(n/2)}
        {\sqrt\pi\,\Gamma((n-1)/2)}
                (1+z_{33}^2)^{-n/2}.
\]
Its upper tail is the source
equation~(117), with its cutoff
chosen to make that tail \(\varepsilon\).
The two conditional cap comparisons
therefore reproduce both source
orientations with their exact scale.
\end{proof}
\end{annexitem}

\begin{annexitem}{annex:p15.student.worked.cases}{thm:p15-student-worked-cases}{Finite-sample calculations and failed premises}
\begin{proof}
For \(\nu=1\), substitution in the
normalized Student density and
\(\Gamma(1)=1,\Gamma(1/2)=\sqrt\pi\)
give \(f_1(t)=1/(\pi(1+t^2))\).
Integration gives its CDF and inversion
gives the tangent quantile.
For \((1,3)\), the mean is \(2\),
the two residual squares sum to \(2\),
and \(e=\sqrt{2/(2\cdot1)}=1\).
Since \(t_{1,3/4}=1\), the two
one-sided intervals and the central
coefficient-\(1/2\) interval follow.
For central coefficient \(3/4\)
the quantile probability is \(7/8\),
and
\(t_{1,7/8}=\tan(3\pi/8)=1+\sqrt2\).
Adding and subtracting this from \(2\)
gives \([1-\sqrt2,3+\sqrt2]\).
These are procedure coverage
coefficients, not probabilities
assigned to the now-fixed observed
parameter intervals.

For \(\nu=2\), the gamma recurrence
gives \(\Gamma(3/2)=\sqrt\pi/2\);
the density is
\[
 \frac{1}{2\sqrt2}(1+t^2/2)^{-3/2}
                       =(2+t^2)^{-3/2}.
\]
The derivative of \(t/(2\sqrt{2+t^2})\)
is this density. Its limits at the two
infinities are \(1/2\) and \(-1/2\),
giving the stated CDF. Solving
\(r=t/\sqrt{2+t^2}\), \(r=2p-1\),
and retaining the sign gives
\(t=\sqrt2r/\sqrt{1-r^2}\).
For \((0,1,2)\), the mean is \(1\),
\(V=2\), and \(e=1/\sqrt3\).
At central coefficient \(9/10\),
\(p=19/20\), \(r=9/10\), and
\(t_{2,19/20}=9\sqrt{2/19}\).
Multiplication by \(e\) gives
\(3\sqrt{6/19}\), as claimed.

For \(\nu=1\), the positive first
moment integral contains
\(\int_1^\infty t/(1+t^2)\,dt=\infty\).
The negative part also diverges by
symmetry; a symmetric density alone
does not define a finite mean.
When \(n=1\), \(V=0\) identically and
\(\nu=0\), so neither \(S^2=V/\nu\)
nor the proved positive-shape gamma
construction exists. This theorem
has expressly assumed \(n\geq2\).

For the nonnormal boundary example,
let \(X_1,X_2\) be independent fair
signs. With probability \(1/2\) they
are equal, in which case \(\bar X\)
is \(1\) or \(-1\) and \(V=0\).
With probability \(1/2\) they differ,
in which case \(\bar X=0\) and \(V=2\).
Thus
\[
 P(\bar X=0,V=2)=\tfrac12
       \ne P(\bar X=0)P(V=2)=\tfrac14.
\]
This verifies dependence directly.
Finally, under the normal model with
\(n=2\), exact pivot coverage gives
\[
 P_{\mu,\sigma}(|\bar X-\mu|\leq ce)
                 =F_1(c)-F_1(-c)
                 =\frac{2\arctan(c)}{\pi}.
\]
For example \(c=1\) gives \(1/2\).
The normal interval with the same
numerical cutoff has central normal
probability \(\int_{-1}^1
(2\pi)^{-1/2}e^{-t^2/2}\,dt>1/2\):
using \(e^{-t^2/2}\geq1-t^2/2\)
on this interval gives the lower bound
\(5/(3\sqrt{2\pi})>1/2\).
This makes the coverage change a
proved finite-sample fact.
\end{proof}
\end{annexitem}
